% GENERATED FILE. Edit canonical lessons, then run npm run generate:books.
% canonical-source-hash: fnv1a64:c6755de538a588eb
% canonical-lessons: JA-C45-neko, JA-C45-uma, JA-C45-ushi, JA-C45-tori, JA-C45-kuma

\chapter{Cat, Horse, Cow, Bird, Bear}
\label{ch:ja-cat-horse-cow-bird-bear}
\hlchaptermodality{45}

\begin{chapteropening}
\textbf{By the end of this chapter:} \emph{I can say a cat, a horse, a cow, a bird and a bear.}

Complete the last lesson of chapter 45: i can say a cat, a horse, a cow, a bird and a bear.
\end{chapteropening}

\section[neko]{\ja{ねこ} (neko) — a cat}

\label{lesson:JA-C45-neko}



\begin{quote}
Before the new one: say the Japanese for a seed, then the Japanese for rain; candy.
\end{quote}

\subsection*{You'll want to know: \ja{ねこ}}
\textbf{\ja{ねこ}} — \emph{neko} — \textquotedblleft{}a cat\textquotedblright{}.

\begin{grammarlens}[title={What you've built}]
One more word for this chapter.
\end{grammarlens}

\subsection*{Your turn}
\begin{itemize}
\raggedright
  \item \emph{Say it:} \emph{neko}
  \item \emph{Say it:} \emph{neko}, once more
  \item \emph{Say it:} say \emph{neko}
  \item \emph{Recall:} say the Japanese for a seed, then the Japanese for rain; candy, then say \emph{neko} again
\end{itemize}

\begin{tcolorbox}[breakable,colback=teal!4,colframe=teal!35!black,title={Before you move on}]
What does \ja{ねこ} mean? (\textquotedblleft{}A cat\textquotedblright{}.) Say it once more.
\end{tcolorbox}
\section[uma]{\ja{うま} (uma) — a horse}

\label{lesson:JA-C45-uma}



\begin{quote}
Before the new one: say the Japanese for rain; candy, then the Japanese for a cat.
\end{quote}

\subsection*{You'll want to know: \ja{うま}}
\textbf{\ja{うま}} — \emph{uma} — \textquotedblleft{}a horse\textquotedblright{}.

\begin{grammarlens}[title={What you've built}]
One more word for this chapter.
\end{grammarlens}

\subsection*{Your turn}
\begin{itemize}
\raggedright
  \item \emph{Say it:} \emph{uma}
  \item \emph{Say it:} \emph{uma}, once more
  \item \emph{Say it:} say \emph{uma}
  \item \emph{Recall:} say the Japanese for rain; candy, then the Japanese for a cat, then say \emph{uma} again
\end{itemize}

\begin{tcolorbox}[breakable,colback=teal!4,colframe=teal!35!black,title={Before you move on}]
What does \ja{うま} mean? (\textquotedblleft{}A horse\textquotedblright{}.) Say it once more.
\end{tcolorbox}
\section[ushi]{\ja{うし} (ushi) — a cow}

\label{lesson:JA-C45-ushi}



\begin{quote}
Before the new one: say the Japanese for a cat, then the Japanese for a horse.
\end{quote}

\subsection*{You'll want to know: \ja{うし}}
\textbf{\ja{うし}} — \emph{ushi} — \textquotedblleft{}a cow\textquotedblright{}.

\begin{grammarlens}[title={What you've built}]
One more word for this chapter.
\end{grammarlens}

\subsection*{Your turn}
\begin{itemize}
\raggedright
  \item \emph{Say it:} \emph{ushi}
  \item \emph{Say it:} \emph{ushi}, once more
  \item \emph{Say it:} say \emph{ushi}
  \item \emph{Recall:} say the Japanese for a cat, then the Japanese for a horse, then say \emph{ushi} again
\end{itemize}

\begin{tcolorbox}[breakable,colback=teal!4,colframe=teal!35!black,title={Before you move on}]
What does \ja{うし} mean? (\textquotedblleft{}A cow\textquotedblright{}.) Say it once more.
\end{tcolorbox}
\section[tori]{\ja{とり} (tori) — a bird}

\label{lesson:JA-C45-tori}



\begin{quote}
Before the new one: say the Japanese for a horse, then the Japanese for a cow.
\end{quote}

\subsection*{You'll want to know: \ja{とり}}
\textbf{\ja{とり}} — \emph{tori} — \textquotedblleft{}a bird\textquotedblright{}.

\begin{grammarlens}[title={What you've built}]
One more word for this chapter.
\end{grammarlens}

\subsection*{Your turn}
\begin{itemize}
\raggedright
  \item \emph{Say it:} \emph{tori}
  \item \emph{Say it:} \emph{tori}, once more
  \item \emph{Say it:} say \emph{tori}
  \item \emph{Recall:} say the Japanese for a horse, then the Japanese for a cow, then say \emph{tori} again
\end{itemize}

\begin{tcolorbox}[breakable,colback=teal!4,colframe=teal!35!black,title={Before you move on}]
What does \ja{とり} mean? (\textquotedblleft{}A bird\textquotedblright{}.) Say it once more.
\end{tcolorbox}
\section[kuma]{\ja{くま} (kuma) — a bear}

\label{lesson:JA-C45-kuma}



\begin{quote}
Before the new one: say the Japanese for a cow, then the Japanese for a bird.
\end{quote}

\subsection*{You'll want to know: \ja{くま}}
\textbf{\ja{くま}} — \emph{kuma} — \textquotedblleft{}a bear\textquotedblright{}.

\begin{grammarlens}[title={What you've built}]
One more word for this chapter.
\end{grammarlens}

\subsection*{Your turn}
\begin{itemize}
\raggedright
  \item \emph{Say it:} \emph{kuma}
  \item \emph{Say it:} \emph{kuma}, once more
  \item \emph{Say it:} say \emph{kuma}
  \item \emph{Recall:} say the Japanese for a cow, then the Japanese for a bird, then say \emph{kuma} again
\end{itemize}

\begin{tcolorbox}[breakable,colback=teal!4,colframe=teal!35!black,title={Before you move on}]
What does \ja{くま} mean? (\textquotedblleft{}A bear\textquotedblright{}.) Say it once more.
\end{tcolorbox}
